%% kapa: exam=1 round=26 year=2015 period=2 form=A subject=appraisal_law subject_ko=부동산관계법규 seq=28 number=28
%% source: Q-Net 감정평가사 기출문제 (cst00309) artlSeq=5138001 file=extracted/1008/5138001/제26회_감정평가사 1차_기출/제26회_감정평가사 1차_기출/2교시_감평_A형.pdf
%% conversion: kordoc 4.24.51 -> tools/gen1.py (auto); review=proofread-sonnet; flags=none
\begin{question}{28}
국유재산법령상 국유재산에 관한 설명으로 옳은 것은?
\begin{choices}
  \item 대통령 관저는 공용재산이다.
  \item 기업용재산은 행정재산이 아니고 일반재산이다.
  \item 정부기업이 직접 사업용으로 사용하는 재산은 보존용재산이다.
  \item 총괄청은 일반재산을 공용재산으로 전환하여 관리할 수 있다.
  \item 행정재산은 시효취득의 대상이 된다.
\end{choices}
\end{question}
\answer{1}
